% ============================================================================
% APPENDIX A: SOURCE CODE
% ============================================================================

\chapter{SOURCE CODE}

This appendix contains key source code excerpts from the EEG-based depression detection system implementation. The complete source code is available in the project repository.

\section{Transformer Encoder (transformer\_encoder.py)}

The Transformer encoder processes CWT scalograms using Vision Transformer-style patch embedding and self-attention.

\begin{lstlisting}[language=Python, caption={Patch Embedding Module}]
class PatchEmbedding(nn.Module):
    """Convert scalogram into patch embeddings."""

    def __init__(self, img_size=(64, 128), patch_size=(8, 16),
                 in_channels=1, embed_dim=128):
        super().__init__()
        self.img_size = img_size
        self.patch_size = patch_size
        self.n_patches = (img_size[0] // patch_size[0]) * \
                         (img_size[1] // patch_size[1])
        self.patch_dim = patch_size[0] * patch_size[1] * in_channels
        self.projection = nn.Linear(self.patch_dim, embed_dim)

    def forward(self, x):
        B, C, H, W = x.shape
        ph, pw = self.patch_size

        # Reshape into patches
        x = x.reshape(B, C, H // ph, ph, W // pw, pw)
        x = x.permute(0, 2, 4, 1, 3, 5).reshape(B, -1, self.patch_dim)

        return self.projection(x)
\end{lstlisting}

\begin{lstlisting}[language=Python, caption={Transformer Encoder Forward Pass}]
def forward(self, x):
    """Process scalogram through transformer.

    Args:
        x: Scalogram tensor of shape (batch, 1, 64, 128)

    Returns:
        CLS token representation of shape (batch, 128)
    """
    B = x.shape[0]

    # Patch embedding: (B, 64 patches, 128 dim)
    x = self.patch_embed(x)

    # Prepend learnable CLS token
    cls_tokens = self.cls_token.expand(B, -1, -1)
    x = torch.cat([cls_tokens, x], dim=1)

    # Add positional encoding
    x = x + self.pos_encoding[:, :x.size(1)]
    x = self.dropout(x)

    # Transformer encoder layers
    x = self.transformer(x)
    x = self.norm(x)

    # Return CLS token output
    return x[:, 0]
\end{lstlisting}

\section{GNN Encoder (gnn\_encoder.py)}

The Graph Attention Network models spatial relationships between EEG electrodes.

\begin{lstlisting}[language=Python, caption={Electrode Graph Construction}]
class ElectrodeGraph:
    """Constructs graph structure for EEG electrodes."""

    # Standard 10-20 electrode positions (normalized 3D coordinates)
    ELECTRODE_POSITIONS_10_20 = {
        'Fp1': [-0.31, 0.95, -0.03], 'Fp2': [0.31, 0.95, -0.03],
        'F7': [-0.81, 0.59, -0.03],  'F3': [-0.55, 0.67, 0.50],
        'Fz': [0.00, 0.71, 0.70],    'F4': [0.55, 0.67, 0.50],
        'F8': [0.81, 0.59, -0.03],   'T3': [-1.00, 0.00, -0.03],
        'C3': [-0.71, 0.00, 0.71],   'Cz': [0.00, 0.00, 1.00],
        'C4': [0.71, 0.00, 0.71],    'T4': [1.00, 0.00, -0.03],
        'T5': [-0.81, -0.59, -0.03], 'P3': [-0.55, -0.67, 0.50],
        'Pz': [0.00, -0.71, 0.70],   'P4': [0.55, -0.67, 0.50],
        'T6': [0.81, -0.59, -0.03],  'O1': [-0.31, -0.95, -0.03],
        'O2': [0.31, -0.95, -0.03]
    }

    def _compute_spatial_adjacency(self):
        """Compute k-nearest neighbors adjacency."""
        adj = kneighbors_graph(
            self.positions, n_neighbors=self.spatial_k,
            mode='connectivity', include_self=False
        ).toarray()
        return np.maximum(adj, adj.T)  # Make symmetric
\end{lstlisting}

\begin{lstlisting}[language=Python, caption={GAT Forward Pass}]
def forward(self, x, edge_index, batch):
    """Process electrode features through GAT layers.

    Args:
        x: Node features (N_total, 576)
        edge_index: Graph edges (2, E)
        batch: Batch assignment for nodes

    Returns:
        Graph-level representation (B, 128)
    """
    # Project input features
    x = F.relu(self.input_proj(x))

    # GAT layers with residual connections
    x = x + F.elu(self.gat1(x, edge_index))
    x = self.norm1(x)

    x = x + F.elu(self.gat2(x, edge_index))
    x = self.norm2(x)

    x = self.gat3(x, edge_index)

    # Global pooling over nodes
    return global_mean_pool(x, batch)
\end{lstlisting}

\section{Attention Fusion (attention\_fusion.py)}

The fusion module combines Transformer and GNN features using cross-attention and gating.

\begin{lstlisting}[language=Python, caption={Cross-Attention Module}]
class CrossAttention(nn.Module):
    """Cross-attention where one modality attends to another."""

    def __init__(self, dim, num_heads=4, dropout=0.1):
        super().__init__()
        self.num_heads = num_heads
        self.head_dim = dim // num_heads
        self.scale = self.head_dim ** -0.5

        self.q_proj = nn.Linear(dim, dim)
        self.k_proj = nn.Linear(dim, dim)
        self.v_proj = nn.Linear(dim, dim)
        self.out_proj = nn.Linear(dim, dim)
        self.dropout = nn.Dropout(dropout)

    def forward(self, query, key_value):
        B, N_q, C = query.shape
        N_kv = key_value.shape[1]

        # Multi-head attention
        q = self.q_proj(query).reshape(B, N_q, self.num_heads,
                                        self.head_dim).permute(0, 2, 1, 3)
        k = self.k_proj(key_value).reshape(B, N_kv, self.num_heads,
                                            self.head_dim).permute(0, 2, 1, 3)
        v = self.v_proj(key_value).reshape(B, N_kv, self.num_heads,
                                            self.head_dim).permute(0, 2, 1, 3)

        attn = (q @ k.transpose(-2, -1)) * self.scale
        attn = F.softmax(attn, dim=-1)
        attn = self.dropout(attn)

        out = (attn @ v).transpose(1, 2).reshape(B, N_q, C)
        return self.out_proj(out)
\end{lstlisting}

\section{WPD Feature Extractor (wpd\_extractor.py)}

The WPD module extracts multi-wavelet features from EEG signals.

\begin{lstlisting}[language=Python, caption={WPD Feature Computation}]
class WPDFeatureExtractor:
    """Multi-wavelet Wavelet Packet Decomposition extractor."""

    def __init__(self, wavelets=['db4', 'sym5', 'coif3'], level=5):
        self.wavelets = wavelets
        self.level = level
        self.n_nodes = 2 ** level  # 32 terminal nodes

    def _compute_node_features(self, coeffs):
        """Compute features from wavelet coefficients."""
        eps = 1e-10
        features = {}

        # Energy
        features['energy'] = np.sum(coeffs ** 2)

        # Shannon Entropy
        p = coeffs ** 2
        p_norm = p / (np.sum(p) + eps)
        p_norm = np.clip(p_norm, eps, 1)
        features['entropy'] = -np.sum(p_norm * np.log(p_norm))

        # Statistical features
        features['mean'] = np.mean(coeffs)
        features['std'] = np.std(coeffs)
        features['skewness'] = skew(coeffs) if len(coeffs) > 2 else 0

        return features

    def extract(self, signal):
        """Extract WPD features from signal."""
        all_features = []
        for wavelet in self.wavelets:
            wp = pywt.WaveletPacket(signal, wavelet, maxlevel=self.level)
            for node in wp.get_level(self.level, 'natural'):
                features = self._compute_node_features(node.data)
                all_features.extend(features.values())
        return np.array(all_features)
\end{lstlisting}

\section{Configuration Files}

\begin{lstlisting}[language=Python, caption={Model Configuration (model\_config.yaml)}]
# Model Architecture Configuration
model:
  transformer:
    img_size: [64, 128]
    patch_size: [8, 16]
    d_model: 128
    nhead: 4
    num_layers: 4
    dim_ff: 512
    dropout: 0.1

  gnn:
    in_features: 576
    hidden_dim: 256
    out_dim: 128
    num_heads: 4
    num_layers: 3
    dropout: 0.3
    spatial_k: 6
    functional_threshold: 0.5

  fusion:
    dim: 128
    num_heads: 4
    dropout: 0.1

  classifier:
    hidden_dims: [64, 32]
    dropout: [0.5, 0.3]
\end{lstlisting}

\begin{lstlisting}[language=Python, caption={Training Configuration (training\_config.yaml)}]
# Training Configuration
training:
  cross_validation: loso  # Leave-One-Subject-Out
  epochs: 30
  batch_size: 16
  gradient_accumulation: 2

  optimizer:
    type: adamw
    lr: 0.0001
    weight_decay: 0.01
    betas: [0.9, 0.999]

  scheduler:
    type: cosine_annealing
    T_0: 10
    T_mult: 2
    eta_min: 0.000001

  early_stopping:
    patience: 10
    min_delta: 0.001

  mixed_precision: true
  gradient_clip: 1.0
\end{lstlisting}
